% !PRINTED XII
\weberpage{XII}
\inhaltabschnitt{Zehnter Abschnitt.}{Algebraische Zahlen und Formen.}
\inhaltpar{87}{Ideale und Formen in algebraischen Körpern}{330}
\inhaltpar{88}{Idealklassen und Formenklassen}{333}
\inhaltpar{89}{Komposition der Formen und Multiplikation der Ideale}{335}
\inhaltabschnitt{Elfter Abschnitt.}{Ideale in quadratischen Körpern.}
\inhaltpar{90}{Diskriminante des quadratischen Körpers}{338}
\inhaltpar{91}{Ideale und Formen in quadratischen Körpern}{340}
\inhaltpar{92}{Primideale im quadratischen Körper}{342}
\inhaltpar{93}{Darstellung von Zahlen als Idealnormen}{344}
\inhaltpar{94}{Das quadratische Reziprozitätsgesetz}{345}
\inhaltpar{95}{Äquivalente Formen und Ideale im quadratischen Körper}{347}
\inhaltabschnitt{Zwölfter Abschnitt.}{Ordnungen im quadratischen Körper.}
\inhaltpar{96}{Diskriminanten der Ordnungen}{351}
\inhaltpar{97}{Ordnungen und Ideale}{353}
\inhaltabschnitt{Dreizehnter Abschnitt.}{Äquivalenz nach Zahlgruppen.}
\inhaltpar{98}{Zahlgruppen in den Ordnungen}{358}
\inhaltpar{99}{Äquivalenz in den Ordnungen}{361}
\inhaltpar{100}{Idealklassen nach den Ordnungen}{362}
\inhaltabschnitt{Vierzehnter Abschnitt.}{Komposition der Formen und Ideale.}
\inhaltpar{101}{Komposition in den Ordnungen}{368}
\inhaltpar{102}{Komposition der Ordnungen}{373}
\inhaltabschnitt{Fünfzehnter Abschnitt.}{Geschlechter der quadratischen Formen.}
\inhaltpar{103}{Darstellung von Zahlen durch quadratische Formen}{376}
\inhaltpar{104}{Charaktere und Geschlechter der quadratischen Formen}{380}
\inhaltpar{105}{Anwendung des \gesperrt{Legendreschen} Symbols}{385}
\inhaltpar{106}{Die Geschlechter der Idealklassen}{388}
\inhaltpar{107}{Zusammensetzung der Normenrestgruppen}{389}
\inhaltpar{108}{Normenreste der Primzahlpotenzen}{390}
\inhaltpar{109}{Die Geschlechter der Ideale}{395}
